% ---------------------------------------------------------------- title
\begin{titlepage}
\thispagestyle{empty}
\vspace*{38mm}
{\sffamily\bfseries\footnotesize\color{eblue}A BEGINNER-TO-INTERMEDIATE COURSE\par}
\vspace{14mm}
{\fontsize{64}{66}\selectfont\bfseries\color{navy}musrSim\par}
\vspace{5mm}
{\color{eblue}\rule{62mm}{2pt}\par}
\vspace{10mm}
{\fontsize{30}{36}\selectfont\color{navy}A Hands-On Geant4 Tutorial\par}
\vspace{14mm}
{\large\itshape\color{ink}Simulating muon transport, detection and decay ---\\
from a text macro to a physics result.\par}
\vfill
{\sffamily\bfseries\color{navy}TDLI Muon Physics Group\par}
\vspace{1.5mm}
{\sffamily\small\color{ink}Tsung-Dao Lee Institute \textperiodcentered\ Shanghai Jiao Tong University\par}
\vspace{2.5mm}
{\sffamily\scriptsize\color{muted}Second edition, October 2026 \textperiodcentered\ for \texttt{musrSim-upgrade} \textperiodcentered\ Geant4 11 + ROOT 6 + Python\par}
\end{titlepage}

% ---------------------------------------------------------------- about
\chapter*{About this edition}
\addcontentsline{toc}{chapter}{About this edition}
\markboth{About this edition}{}

This tutorial teaches Geant4 simulation through \musrSim, using the command syntax, ROOT output and build system of the \texttt{musrSim-upgrade} repository (\url{github.com/kimsiang/musrSim-upgrade}). The commands, branch names and file names in this book were checked against that source. The two worked examples --- \code{first_sim.mac} and the capstone \code{tf_musr.mac} --- ship in the repository's \code{run/} folder, and every number quoted from them was produced by actually running them.

\begin{checkbuild}
\musrSim\ is a living, forked codebase --- PSI, SJTU and individual groups each carry local changes. The \emph{standard Geant4} commands (\code{/run}, \code{/vis}, \code{/tracking}, \code{/control}, \code{/gps}) are universal, but the \code{/musr/...} and \code{/gun/...} commands are defined \emph{in your build's source}. If a command in this book is rejected, that is your cue to check the code, not a typo to fight. Chapter~\ref{ch:macro} shows how to list what your executable understands, and Chapter~\ref{ch:trouble} how to find any keyword in the source in ten seconds.
\end{checkbuild}

\eyebrow{What changed since the first edition}
The first edition was written for the older \texttt{musrsim-sms} distribution on Geant4~10.7. This edition follows the cleaned-up \texttt{musrSim-upgrade} code:

\begin{itemize}
  \item \textbf{Geant4 11 and modern CMake.} The code now builds against Geant4~11 (tested with 11.4.3) and ROOT~6 (tested with 6.40), with CMake~$\ge$\,3.16. Geant4~10.7 still works. Field-accuracy limits and optical-property names that changed in Geant4~11 are handled automatically.
  \item \textbf{One physics list, configured from the macro.} \musrSim\ always uses its own \code{musrPhysicsList}, and \emph{your macro's} \code{/musr/command process} lines decide which electromagnetic processes exist. A macro without them simulates muons that never stop (Section~\ref{sec:physicsblock}).
  \item \textbf{Cosmic-ray generators built in.} EcoMug and CRY are now part of the code, driven by \code{/gun/ecomug/...} and \code{/gun/cry/...} commands (Appendix~\ref{app:cosmics}).
  \item \textbf{Run modes.} The executable always needs a macro. Use \code{../musrSim macro.mac idle} to get an interactive prompt after the macro has run.
  \item \textbf{Correct batch seeding.} Only \code{/musr/run/randomOption 1} makes parallel jobs statistically independent (Section~\ref{sec:scale}).
  \item \textbf{Output names} follow \code{data/musrSim_jp<NNN>[_<name>].root}, with a zero-padded run number.
\end{itemize}

\vfill
{\sffamily\scriptsize\color{muted}Typeset with \texttt{xelatex} in Latin Modern (body), TeX Gyre Heros (headings) and DejaVu Sans Mono (code). Source: \texttt{doc/musrSim\_tutorial.tex}; build with \texttt{latexmk -xelatex}. Figures are regenerated by \texttt{doc/scripts/make\_figures.py}.\par}

% ---------------------------------------------------------------- how to use
\chapter*{How to use this book}
\addcontentsline{toc}{chapter}{How to use this book}
\markboth{How to use this book}{}

\textbf{Who this is for.} Anyone who wants to run muon simulations and does not yet know Geant4. We assume you are comfortable in a Unix shell and can read a little C++ and Python. No Geant4, ROOT or Monte-Carlo experience is required.

\textbf{The three things you actually need.} A working \musrSim\ executable (Chapter~\ref{ch:install}), a text editor, and a Python environment for analysis (Chapter~\ref{ch:python}). That is the whole toolchain.

\eyebrow{The learning path}
\begin{center}
\begin{tikzpicture}[node distance=5mm and 5mm]
  \node[box] (a) {\textbf{1--2}\\Orient \& install};
  \node[box,right=of a] (b) {\textbf{3}\\Macro language};
  \node[hot,right=of b] (c) {\textbf{4}\\Geometry};
  \node[hot,right=of c] (d) {\textbf{5}\\Muon beam};
  \node[hot,below=7mm of d] (e) {\textbf{6}\\Hits \& ROOT};
  \node[box,left=of e] (f) {\textbf{7}\\Python analysis};
  \node[box,left=of f] (g) {\textbf{8}\\Fields, physics,\\scale};
  \node[box,left=of g] (h) {\textbf{9--10}\\Capstone \&\\troubleshooting};
  \draw[arr] (a)--(b); \draw[arr] (b)--(c); \draw[arr] (c)--(d); \draw[arr] (d)--(e);
  \draw[arr] (e)--(f); \draw[arr] (f)--(g); \draw[arr] (g)--(h);
\end{tikzpicture}
\end{center}
{\small\color{muted}Shaded = the core craft. If you only have an afternoon, do Chapters 2, 3, 4 and 6.\par}

\eyebrow{How to read the callouts}
\begin{keypoint}The one idea on the page that, if you remember nothing else, you should remember.\end{keypoint}
\begin{tip}A shortcut, a good habit, or a small trick that saves you time.\end{tip}
\begin{pitfall}A mistake beginners actually make. Read these --- they are the errors that cost an afternoon.\end{pitfall}
\begin{physics}The physics behind the code: why the numbers come out the way they do.\end{physics}
\begin{hood}What Geant4 or \musrSim\ is doing internally. Optional on a first read, useful when debugging.\end{hood}
\begin{checkbuild}Something that may differ in your fork. A prompt to verify against your own source, not to trust blindly.\end{checkbuild}

\eyebrow{Conventions for code}
Shell commands, macro files (\code{.mac}), C++ and Python are shown in monospaced blocks with a blue rule on the left and a caption underneath. A long line that does not fit is wrapped and marked with \(\hookrightarrow\); type it as \emph{one} line. Comments begin with \code{#}. Paths such as \code{build/run} are relative to the top of the repository.

\begin{pitfall}
The single most important convention in this book. In every \code{/musr/command construct} line, box and tube dimensions are \emph{half-lengths}, following Geant4's \code{G4Box} and \code{G4Tubs}. A ``\code{100 100 5}'' box is \(200\times200\times10\)\,mm. Beginners routinely build detectors twice the size they intended. We flag this again in Chapter~\ref{ch:geometry}, but internalise it now.
\end{pitfall}

{\hypersetup{linkcolor=ink}\tableofcontents}
